% concatenated from QEC_and_EAQEC_Codes_from_Hermitian_Sums_and_Hulls_of_Cyclic_Codes_over_F_2x_F_2_vF_2_.tex
%Version 3.1 December 2024
% See section 11 of the User Manual for version history
%
%%%%%%%%%%%%%%%%%%%%%%%%%%%%%%%%%%%%%%%%%%%%%%%%%%%%%%%%%%%%%%%%%%%%%%
%%                                                                 %%
%% Please do not use \input{...} to include other tex files.       %%
%% Submit your LaTeX manuscript as one .tex document.              %%
%%                                                                 %%
%% All additional figures and files should be attached             %%
%% separately and not embedded in the \TeX\ document itself.       %%
%%                                                                 %%
%%%%%%%%%%%%%%%%%%%%%%%%%%%%%%%%%%%%%%%%%%%%%%%%%%%%%%%%%%%%%%%%%%%%%

%%\documentclass[referee,sn-basic]{sn-jnl}% referee option is meant for double line spacing

%%=======================================================%%
%% to print line numbers in the margin use lineno option %%
%%=======================================================%%

%%\documentclass[lineno,pdflatex,sn-basic]{sn-jnl}% Basic Springer Nature Reference Style/Chemistry Reference Style

%%=========================================================================================%%
%% the documentclass is set to pdflatex as default. You can delete it if not appropriate.  %%
%%=========================================================================================%%

%%\documentclass[sn-basic]{sn-jnl}% Basic Springer Nature Reference Style/Chemistry Reference Style

%%Note: the following reference styles support Namedate and Numbered referencing. By default the style follows the most common style. To switch between the options you can add or remove “Numbered” in the optional parenthesis. 
%%The option is available for: sn-basic.bst, sn-chicago.bst%  
 
%%\documentclass[pdflatex,sn-nature]{sn-jnl}% Style for submissions to Nature Portfolio journals
%%\documentclass[pdflatex,sn-basic]{sn-jnl}% Basic Springer Nature Reference Style/Chemistry Reference Style
\documentclass[pdflatex,sn-mathphys-num]{sn-jnl}% Math and Physical Sciences Numbered Reference Style
%%\documentclass[pdflatex,sn-mathphys-ay]{sn-jnl}% Math and Physical Sciences Author Year Reference Style
%%\documentclass[pdflatex,sn-aps]{sn-jnl}% American Physical Society (APS) Reference Style
%%\documentclass[pdflatex,sn-vancouver-num]{sn-jnl}% Vancouver Numbered Reference Style
%%\documentclass[pdflatex,sn-vancouver-ay]{sn-jnl}% Vancouver Author Year Reference Style
%%\documentclass[pdflatex,sn-apa]{sn-jnl}% APA Reference Style
%%\documentclass[pdflatex,sn-chicago]{sn-jnl}% Chicago-based Humanities Reference Style

%%%% Standard Packages
%%<additional latex packages if required can be included here>

\usepackage{graphicx}%
\usepackage{multirow}%
\usepackage{amsmath,amssymb,amsfonts}%
\usepackage{amsthm}%
\usepackage{mathrsfs}%
\usepackage[title]{appendix}%
%\usepackage{xcolor}%
%\usepackage{textcomp}%
%\usepackage{manyfoot}%
%\usepackage{booktabs}%
\usepackage{algorithm}%
\usepackage{algorithmicx}%
\usepackage{algpseudocode}%
\usepackage{listings}%
\usepackage{faktor}
\usepackage{hyperref}
\usepackage{stmaryrd}
%%%%

%%%%%=============================================================================%%%%
%%%%  Remarks: This template is provided to aid authors with the preparation
%%%%  of original research articles intended for submission to journals published 
%%%%  by Springer Nature. The guidance has been prepared in partnership with 
%%%%  production teams to conform to Springer Nature technical requirements. 
%%%%  Editorial and presentation requirements differ among journal portfolios and 
%%%%  research disciplines. You may find sections in this template are irrelevant 
%%%%  to your work and are empowered to omit any such section if allowed by the 
%%%%  journal you intend to submit to. The submission guidelines and policies 
%%%%  of the journal take precedence. A detailed User Manual is available in the 
%%%%  template package for technical guidance.
%%%%%=============================================================================%%%%

%% as per the requirement new theorem styles can be included as shown below
\theoremstyle{thmstyleone}%
\newtheorem{theorem}{Theorem}%  meant for continuous numbers
%%\newtheorem{theorem}{Theorem}[section]% meant for sectionwise numbers
%% optional argument [theorem] produces theorem numbering sequence instead of independent numbers for Proposition
\newtheorem{proposition}[theorem]{Proposition}% 
%%\newtheorem{proposition}{Proposition}% to get separate numbers for theorem and proposition etc.
\newtheorem{lemma}[theorem]{Lemma}
\theoremstyle{thmstyletwo}%
\newtheorem{example}{Example}%
\newtheorem{remark}{Remark}%
\newtheorem{corollary}{Corollary}%

\theoremstyle{thmstylethree}%
\newtheorem{definition}{Definition}%

\raggedbottom
%%\unnumbered% uncomment this for unnumbered level heads

\begin{document}

\title[QEC and EAQEC Codes from Hermitian Sums and Hulls of Cyclic Codes over $\mathbb{F}_2 \times (\mathbb{F}_2+v\mathbb{F}_2)$]{QEC and EAQEC Codes from Hermitian Sums and Hulls of Cyclic Codes over $\mathbb{F}_2 \times (\mathbb{F}_2+v\mathbb{F}_2)$}

%%=============================================================%%
%% GivenName	-> \fnm{Joergen W.}
%% Particle	-> \spfx{van der} -> surname prefix
%% FamilyName	-> \sur{Ploeg}
%% Suffix	-> \sfx{IV}
%% \author*[1,2]{\fnm{Joergen W.} \spfx{van der} \sur{Ploeg} 
%%  \sfx{IV}}\email{iauthor@gmail.com}
%%=============================================================%%

\author[1]{\fnm{Rabia} \sur{Zengin}}\email{rabia.zengin@bilgi.edu.tr}

\author*[2]{\fnm{Mehmet Emin} \sur{Köroğlu}}\email{mkoroglu@yildiz.edu.tr}
\equalcont{These authors contributed equally to this work.}

\affil[1]{\orgdiv{Department of Mathematics}, \orgname{Istanbul Bilgi University}, \city{Istanbul}, \postcode{34060}, \state{Eyüpsultan},  \country{Türkiye}}

\affil*[2]{\orgdiv{Department of Mathematics}, \orgname{Yildiz Technical University}, \city{Istanbul}, \postcode{34220}, \state{Esenler}, \country{Türkiye}}

%%==================================%%
%% Sample for unstructured abstract %%
%%==================================%%

\abstract{In this work, we determine the generator polynomials for the Hermitian hulls and Hermitian sums of cyclic codes defined over the composite ring $\mathbb{F}_2 \times (\mathbb{F}_2 + v\mathbb{F}_2)$, where $v^2 = v$. Based on these structures, we develop quantum error-correcting (QEC) codes by applying the Hermitian dual version of Quantum Construction~X to the obtained Hermitian hulls and sums. Moreover, by employing matrix product code methods on linear complementary dual (LCD) codes defined over the same ring, we derive families of entanglement-assisted quantum error-correcting (EAQEC) codes.}

%%================================%%
%% Sample for structured abstract %%
%%================================%%

% \abstract{\textbf{Purpose:} The abstract serves both as a general introduction to the topic and as a brief, non-technical summary of the main results and their implications. The abstract must not include subheadings (unless expressly permitted in the journal's Instructions to Authors), equations or citations. As a guide the abstract should not exceed 200 words. Most journals do not set a hard limit however authors are advised to check the author instructions for the journal they are submitting to.
% 
% \textbf{Methods:} The abstract serves both as a general introduction to the topic and as a brief, non-technical summary of the main results and their implications. The abstract must not include subheadings (unless expressly permitted in the journal's Instructions to Authors), equations or citations. As a guide the abstract should not exceed 200 words. Most journals do not set a hard limit however authors are advised to check the author instructions for the journal they are submitting to.
% 
% \textbf{Results:} The abstract serves both as a general introduction to the topic and as a brief, non-technical summary of the main results and their implications. The abstract must not include subheadings (unless expressly permitted in the journal's Instructions to Authors), equations or citations. As a guide the abstract should not exceed 200 words. Most journals do not set a hard limit however authors are advised to check the author instructions for the journal they are submitting to.
% 
% \textbf{Conclusion:} The abstract serves both as a general introduction to the topic and as a brief, non-technical summary of the main results and their implications. The abstract must not include subheadings (unless expressly permitted in the journal's Instructions to Authors), equations or citations. As a guide the abstract should not exceed 200 words. Most journals do not set a hard limit however authors are advised to check the author instructions for the journal they are submitting to.}

\keywords{Hermitian hulls of cyclic codes, Hermitian sums of cyclic codes, Quantum codes, Entanglement-assisted quantum codes, Quantum construction X}

%%\pacs[JEL Classification]{D8, H51}

\pacs[MSC Classification]{94B05, 94B15}

\maketitle

\section{Introduction}\label{sec1}

Quantum error correction (QEC) emerged alongside growing interest in quantum computers. In some research \cite{benioff1980computer,deutsch1985quantum,feynman2018simulating,simon1997power,yao1993quantum}, it is suggested that quantum computers could outperform classical ones on certain problems. The computational intractability of integer factorization constitutes the foundation of widely implemented public-key cryptographic schemes, such as RSA, which ensure the security of online communications and banking system. QEC techniques are essential for mitigating decoherence in quantum systems; without them, quantum computers would be limited to trivial problem sizes.

Several constructions for these codes have been developed, among which the Calderbank-Shor-Steane (CSS) construction \cite{calderbank1997quantum,steane1996multiple} is particularly significant, as it derives a relationship between classical and quantum codes. In \cite{kai2012new}, quantum MDS codes are obtained from negacyclic codes. In \cite{grassl1999quantum}, quantum BCH codes are constructed. In \cite{ashraf2021new}, some quantum and LCD codes are obtained.

Hsieh et al. \cite{hsieh2007general} introduced a foundational class of quantum codes, termed as entanglement-assisted quantum error-correcting codes (EAQECCs), which coherently integrate the theoretical benefits of entanglement-assisted and operator-based quantum error correction frameworks. In \cite{cao2025entanglement}, EAQECCs are obtained by means of matrix-product codes. In \cite{pereira2021entanglement}, EAQECCs are constructed from algebraic geometry codes. One can look at the references \cite{li2025eaqec,koroglu2019new,pang2021new,sari2021new,chen2018entanglement,li2019entanglement,pereira2022entanglement,sok2022linear}.

In \cite{euclideansum}, QEC and EAQEC codes are obtained by means of Euclidean sums and hulls of cyclic codes over the ring $\mathbb{F}_2\times(\mathbb{F}_2+v\mathbb{F}_2)$. Inspired by \cite{euclideansum}, we will obtain QEC and EAQEC codes via Hermitian sums and hulls of cyclic code over the ring. Specifically, the codes we obtained in Theorem \ref{th12} using Hermitian sums and hulls are derived differently from the code parameters in Theorem 3.5 in \cite{euclideansum}.


This paper is organized as follows. In Section \ref{sec2}, we discuss the ring $\mathbb{F}_2\times(\mathbb{F}_2+v\mathbb{F}_2)$, along with essential definitions and lemmas related to cyclic codes, the Gray map, Euclidean hulls and sums, and Hermitian hulls and sums. In Section \ref{sec3}, we examine the relationship between Hermitian hulls and sums of a linear code under the Gray map. By considering the framework of cyclic codes over $\mathbb{F}_2\times(\mathbb{F}_2+v\mathbb{F}_2)$, we determine the generators of the Hermitian hull and sum of a cyclic code over the ring. Then, we construct QEC codes from these Hermitian hulls and sums in Section \ref{sec4}. Thanks to Hermitian structure, we constructed new QEC codes that different from the codes obtained by Euclidean case. Furthermore, we obtain EAQEC codes through Quantum Construction X and matrix product codes in Section \ref{sec5}. The last section concludes the paper.


\section{Preliminaries}
\label{sec2}
The finite commutative ring
\begin{equation*}
\mathcal{R}:=\mathbb{F}_2\times(\mathbb{F}_2+v\mathbb{F}_2)=\{(u_1,
u_2+vu_3) : u_i\in\mathbb{F}_2, i\in\{1,2,3\}, v^2=v\}
\end{equation*}
is introduced in \cite{F2ring} and then linear and cyclic codes over $\mathcal{R}$ are constructed in \cite{F2ring,cyclicF2,selfdualcodeinring,eaqecinring}.

\noindent For vectors $\mathbf{t}=(t_{1},\dots ,t_{n})$ and $\mathbf{w}=(w_{1},\dots ,w_{n})\in \mathcal{R}=\mathbb{F}_{2}\times (\mathbb{F}_{2}+v\mathbb{F}_{2})$, we define the \emph{Euclidean and} \emph{Hermitian inner
product} by
\begin{equation*}
\left\langle \mathbf{t},\mathbf{w}\right\rangle _{E}=\sum_{i=1}^{n}t_{i}w_{i}
\text{ and }\left\langle \mathbf{t},\mathbf{w}\right\rangle
_{H}=\sum_{i=1}^{n}t_{i}\overline{w_{i}},
\end{equation*}
respectively, where $\overline{w_{i}}$ denotes the conjugate of the element $w_{i}$ under the conjugation map in $\mathbb{F}_{2}+v\mathbb{F}_{2}$. The
conjugation on $\mathbb{F}_{2}+v\mathbb{F}_{2}$ is given by $\overline{b+vc}=b+c+vc$ and it acts component-wise on $\mathcal{R}$ as $\overline{(a,b+vc)}
=(a,b+c+vc).$

\noindent Suppose that $\mathcal{C}$ is a linear code over $\mathcal{R}$.
Then the Euclidean dual and Hermitian dual code of $\mathcal{C}$ are
\begin{equation*}
\mathcal{C}^{\perp }=\{\mathbf{t}\in \mathcal{R}^{n}:\left\langle \mathbf{t},
\mathbf{w}\right\rangle _{E}=(0,0)\text{ for all }\mathbf{w}\in \mathcal{C}%
\},
\end{equation*}
and
\begin{equation*}
\mathcal{C}^{\perp _{H}}=\{\mathbf{t}\in \mathcal{R}^{n}:\left\langle
\mathbf{t},\mathbf{w}\right\rangle _{H}=(0,0)\text{ for all }\mathbf{w}\in
\mathcal{C}\}.
\end{equation*}
\noindent In \cite{F2ring}, the Gray map $\tau $ from $\mathcal{R}$ to $\mathbb{F}_{2}^{3}$ is defined as
\begin{equation*}
\tau :\mathcal{R}\rightarrow \mathbb{F}_{2}^{3},\text{ \ }
(a_{1},a_{2}+va_{3})\mapsto (a_{1}+a_{2},a_{1}+a_{3},a_{1}+a_{2}+a_{3})
\end{equation*}
for all $a_{1},$ $a_{2},$ $a_{3}\in \mathbb{F}_{2}$. The map extends
naturally to $\mathcal{R}^{n}$ as
\begin{equation*}
\tau :\mathcal{R}^{n}\rightarrow \mathbb{F}_{2}^{3n},\text{ \ }(\mathbf{a},
\mathbf{b}+v\mathbf{c})\mapsto (\mathbf{a+b},\mathbf{a+c},\mathbf{a+b}+
\mathbf{c})
\end{equation*}
for all $\mathbf{a},$ $\mathbf{b},$ $\mathbf{c}\in \mathbb{F}_{2}^{n}$.


Suppose that $\mathcal{C}$ is a linear code of length $n$. If $(c_{n-1}, c_0,\dots , c_{n-2})\in \mathcal{C}$ whenever $(c_0, c_1,\dots, c_{n-1})\in \mathcal{C}$, then $\mathcal{C}$ is termed as a cyclic code.


\noindent The polynomial $l(x)\in\mathcal{R}[x]$ can be considered as $l_0+l_1x+l_2x^2+\dots l_{n-1}x^{n-1}$, where $l_j=(a_j,b_j+vc_j)$ for $j\in\{0,1,\dots ,n-1\}$. Let $a(x)=a_0+a_1x+\dots +a_{n-1}x^{n-1}$, $b(x)=b_0+b_1x+\dots +b_{n-1}x^{n-1}$ and $c(x)=c_0+c_1x+\dots +c_{n-1}x^{n-1}$ be in $\mathbb{F}_2[x]$. Then $l(x)$ can be written as $l(x)=(a(x),b(x)+vc(x))$. The multiplication in $\mathcal{R}[x]$ is defined as
\begin{align*}
l(x)t(x) & = (l_1(x),l_2(x)+vl_3(x))(t_1(x),t_2(x)+vt_3(x)) \\
 & = (l_1(x)t_1(x),l_2(x)t_2(x)+v(l_2(x)t_3(x)+l_3(x)t_2(x)+ l_3(x)t_3(x) )),
\end{align*}
where $l(x), t(x)\in \mathcal{R}[x]$.
Let $\mathbf{1}=(1,1)\in\mathcal{R}$. Further, there is a one-to-one correspondence between the polynomial $t(x)=t_0+t_1x+\dots +t_{n-1}x^{n-1}\in \mathcal{R}[x]/\langle \mathbf{1}x^{n}-\mathbf{1} \rangle$ and the element $t=(t_0,t_1,\dots ,t_{n-1})\in\mathcal{R}^n$.

The Hamming weight $wt(x)$ of an element $x$ over $\mathbb{F}_2$ is the number of its nonzero coordinates. In \cite{F2ring}, the Gray weight of $c$ in $\mathcal{R}$ is defined as
\begin{equation*}
wt_G(c) =
\begin{cases}
0, & \text{if } c = (0,0), \\
1, & \text{if } c = (1,1), (1,v), (1,1+v), \\
2, & \text{if } c = (0,1), (0,v), (0,1+v), \\
3, & \text{if } c = (1,0).
\end{cases}
\end{equation*}
Then the Gray weight of $\mathbf{c}=(c_1,c_2,\dots ,c_n)$ in $\mathcal{R}^n$ is $wt_G(\mathbf{c})=\sum_{i=1}^{n}wt_G(c_i)$. The smallest nonzero Gray weight among all codewords of a code is called the Gray weight
of the code. Let $\mathbf{c},\mathbf{c}'\in\mathcal{R}^n$. The Gray
distance between $\mathbf{c}$ and $\mathbf{c}'$ is $d_G(\mathbf{c},\mathbf{c}')=wt_G(\mathbf{c}-\mathbf{c}')$. The smallest nonzero Gray distance between all pairs of distinct
codewords of a code is termed as the
Gray distance of the code.

\begin{lemma} (\cite{F2ring}, Proposition 4, Theorem 3)
\label{graymap} The Gray map $\tau$ defined above satisfies the followings.

\begin{itemize}
\item[(i)] It is an orthogonality preserving map.
\item[(ii)] It is a weight preserving map from $\mathcal{R}^n$ to $\mathbb{F}_{2}^{3n}$ (i.e., from Gray weight to Hamming weight).
\end{itemize}
\end{lemma}

In the following lemma, we recall from Corollary 3 in \cite{F2ring} that the relation between the Euclidean dual of the image of a code over the ring $\mathcal{R}$ under the Gray map $\tau$ and the image of the Hermitian dual of the code under the Gray map $\tau$.

\begin{lemma} (\cite{F2ring}, Corollary 3)
\label{hull} Assume that $\mathcal{C}^{\perp _{H}}$ is the Hermitian dual of
the code $\mathcal{C}$ over $\mathcal{R}$. Then $\tau (\mathcal{C}^{\perp
_{H}})=\tau (\mathcal{C})^{\perp }$.
\end{lemma}


The Euclidean and Hermitian hulls of a linear code $\mathcal{C}$ over the ring $\mathcal{R}$ is defined as $\mathcal{C} \cap \mathcal{C}^{\perp}$ and $\mathcal{C} \cap \mathcal{C}^{\perp_H}$ and denoted by $\mathrm{Hull_E}(\mathcal{C})$ and $\mathrm{Hull_H}(\mathcal{C})$, respectively. If $\mathcal{C}_1$ and $\mathcal{C}_2$ are codes of length $n$ over $\mathcal{R}$, then $\mathcal{C}_1+\mathcal{C}_2= \{ \mathbf{c'} +
\mathbf{c''} : \mathbf{c'}\in \mathcal{C}_1,
\mathbf{c''} \in \mathcal{C}_2 \}$  is the sum of $\mathcal{C}_1$ and $\mathcal{C}_2$. In particular, we denote $\mathcal{C}+\mathcal{C}^{\perp}$ by $\mathrm{Sum_E}(\mathcal{C})$ and $\mathcal{C}+\mathcal{C}^{\perp_H}$ by $\mathrm{Sum_H}(\mathcal{C})$, which are termed as the Euclidean and Hermitian sum of $\mathcal{C}$, respectively.


\section{Hermitian Hulls and Sums of Linear Codes over $\mathbb{F}_2 \times (\mathbb{F}_2+v\mathbb{F}_2)$}
\label{sec3}
In this section, we obtain the Hermitian hulls and sums of cyclic codes over $\mathcal{R}$. First of all, we determine the relation between the Euclidean and Hermitian sums and hulls of linear codes over $\mathcal{R}$.

\begin{lemma}\label{l2}
Suppose that $\mathcal{C}$ is a linear code over $\mathcal{R}$. Then
\begin{itemize}
    \item[(i)] $\tau(\mathrm{Hull_H}(\mathcal{C}))=\mathrm{Hull_E}(\tau(\mathcal{C}))$.
    \item[(ii)] $\mathrm{Sum_E}(\tau(\mathcal{C}))=\tau(\mathrm{Sum_H}(\mathcal{C}))$.
\end{itemize}
\end{lemma}

\begin{proof}

\begin{itemize}
\item[(i)] Let $\mathbf{a}\in \tau (\mathrm{Hull_{H}}(\mathcal{C}))$. Since $%
\tau $ is onto, there exists $\mathbf{x}\in \mathrm{Hull_{H}}(\mathcal{C})$
such that $\tau (\mathbf{x})=\mathbf{a}$. Then $\mathbf{x}\in \mathcal{C}\cap \mathcal{C}^{\perp _{H}}$ and therefore $\mathbf{x}\in \mathcal{C}$ and
$\mathbf{x}\in \mathcal{C}^{\perp _{H}}$. So, $\tau (\mathbf{x})\in \tau (\mathcal{C})$ and $\tau (\mathbf{x})\in \tau (\mathcal{C}^{\perp _{H}})=\tau
(\mathcal{C})^{\perp }$ by Lemma \ref{hull}. Then $\mathbf{a}=\tau (\mathbf{x})\in \tau (\mathcal{C})\cap \tau (\mathcal{C})^{\perp }=\mathrm{Hull_{E}}(\tau (\mathcal{C}))$. Thus, $\tau (\mathrm{Hull_{H}}(\mathcal{C}))\subseteq
\mathrm{Hull_{E}}(\tau (\mathcal{C}))$. On the other hand, assume $\mathbf{a}\in \mathrm{Hull_{E}}(\tau (\mathcal{C}))$. By Lemma \ref{hull}, $\mathbf{a}\in \tau (\mathcal{C})$ and $\mathbf{a}\in \tau (\mathcal{C})^{\perp }=\tau (\mathcal{C}^{\perp _{H}})$. So, there exists some $\mathbf{x}\in \mathcal{C}$
and $\mathbf{y}\in \mathcal{C}^{\perp _{H}}$ such that $\tau (\mathbf{x})=\mathbf{a}$ and $\tau (\mathbf{y})=\mathbf{a}$. Since $\tau $ is injective, $\mathbf{x}=\mathbf{y}$ and therefore $\mathbf{x}\in \mathcal{C}\cap \mathcal{C}^{\perp _{H}}=\mathrm{Hull_{H}}(\mathcal{C})$. Then $\mathbf{a}=\tau
(x)\in \tau (\mathrm{Hull_{H}}(\mathcal{C}))$. Thus, $\tau (\mathrm{Hull_{H}}(\mathcal{C}))\supseteq \mathrm{Hull_{E}}(\tau (\mathcal{C}))$ and hence $\tau (\mathrm{Hull_{H}}(\mathcal{C}))=\mathrm{Hull_{E}}(\tau (\mathcal{C}))$.

\item[(ii)] Let $\mathbf{a}\in \tau (\mathrm{Sum_{H}}(\mathcal{C}))=\tau (\mathcal{C}+\mathcal{C}^{\perp _{H}})$. Then there exists $\mathbf{x}\in \mathcal{C}$ and $\mathbf{y}\in \mathcal{C}^{\perp _{H}}$ such that $\mathbf{a}=\tau (\mathbf{x}+\mathbf{y})$. Since $\tau $ is linear, $\mathbf{a}=\tau (\mathbf{x})+\tau (\mathbf{y})$. Also, $\tau (\mathbf{x})\in \tau (\mathcal{C})$ and $\tau (\mathbf{y})\in \tau (\mathcal{C}^{\perp _{H}})=\tau (\mathcal{C})^{\perp }$. Then $\mathbf{a}\in \tau (\mathcal{C})+\tau (\mathcal{C})^{\perp }=\mathrm{Sum_{E}}(\tau (\mathcal{C}))$. So, $\tau (\mathrm{Sum_{H}}(\mathcal{C}))\subseteq \mathrm{Sum_{E}}(\tau (\mathcal{C}))$. Now, let $\mathbf{b}\in \mathrm{Sum_{E}}(\mathcal{C})$. Then $\mathbf{b}\in \tau (\mathcal{C})+\tau (\mathcal{C}^{\perp _{H}})$ and therefore $\mathbf{b}=\tau (\mathbf{u})+\tau (\mathbf{v})$ for some $\mathbf{u}\in \mathcal{C}$ and $\mathbf{v}\in \mathcal{C}^{\perp _{H}}$. So, $\mathbf{b}=\tau (\mathbf{u}+\mathbf{v})\in \tau (\mathrm{Sum_{H}}(\mathcal{C}))$. So, $\tau (\mathrm{Sum_{H}}(\mathcal{C}))\supseteq \mathrm{Sum_{E}}(\tau (\mathcal{C}))$. Thus, $\tau (\mathrm{Sum_{H}}(\mathcal{C}))=\mathrm{Sum_{E}}(\tau (\mathcal{C}))$.
\end{itemize}
\end{proof}

\noindent We denote the monic reciprocal polynomial of $h(x)=t_0+t_1x+\dots +t_nx^n$ by $\Bar{h}(x)$, i.e., $\Bar{h}(x)=t_n^{-1}x^{n}h(\frac{1}{x})$.

\begin{theorem}\cite{cyclicF2}\label{7.1}
Let $\mathcal{C}$ be a linear code over $\mathcal{R}$. Then
\begin{itemize}
\item[(i)] $\mathcal{C}=\langle \nu(x)\rangle$, where $\nu(x)=(r_1(x),(1+v)r_2(x)+r_3(x))$ and $r_i(x)\mid (x^n-1)$ for all $i\in \{1,2,3\}$. Also, the polynomial $\nu(x)$ is unique and $|\mathcal{C}|=2^{3n-\mathrm{deg}(r_1(x))-\mathrm{deg}(r_2(x))-\mathrm{deg}(r_3(x))}$.
\item[(ii)] $\mathcal{C}^{\perp_H}=\langle \Bar{h}(x)\rangle$, where $\Bar{h}(x)=(\Bar{h}_1(x),v\Bar{h}_2(x)+(1+v)\Bar{h}_3(x))$ and $x^n-1=r_i(x)h_i(x)$ for all $i\in \{1,2,3\}$. Further, $|\mathcal{C}^{\perp_H}|=2^{\mathrm{deg}(r_1(x))+\mathrm{deg}(r_2(x))+\mathrm{deg}(r_3(x))}$.
\end{itemize}
\end{theorem}

\begin{remark}\cite{F2ring}\label{remark7.1}
A linear code $\mathcal{C}$ over $\mathcal{R}$ of length $n$ is permutation equivalent to direct product of $\mathcal{C}_1 $ and $ \mathcal{C}_2$, where $\mathcal{C}_1$ is a binary linear code of length $n$, $\mathcal{C}_2$ is a linear code over $\mathcal{R}$ of length $n$, which will be denoted by $\mathcal{C}= (\mathcal{C}_1,\mathcal{C}_2)$.
\end{remark}

\begin{theorem}\cite{cyclicinring}\label{th7.2}
Suppose that $\Theta$ is the Gray map from $\mathbb{F}_2+v\mathbb{F}_2$ to $\mathbb{F}_2^2$ defined by $\Theta(x+vy)=(x,x+y)$. Let $\mathcal{C}$ be a linear code of length $n$ over $\mathbb{F}_2+v\mathbb{F}_2$.
Then $\Theta(\mathcal{C})=\mathcal{C}_1 \otimes \mathcal{C}_2:=\{(x,y):x\in\mathcal{C}_1,y\in\mathcal{C}_2)\}$ and $|\mathcal{C}|=|\mathcal{C}_1||\mathcal{C}_2|$, where $\mathcal{C}_1=\{x\in\mathbb{F}_2^n : x+vy \in\mathcal{C}\}$ and $\mathcal{C}_2=\{x+y\in\mathbb{F}_2^n : x+vy \in\mathcal{C}\}$. Also, $\Theta(\mathcal{C})$ is linear.
\end{theorem}

\noindent By utilizing Theorem \ref{th7.2}, we get the following proposition.

\begin{proposition}\label{prop1}
Let $d$ and $d_G$ be the minimum Hamming and Gray distances of a linear code $\mathcal{C}$ over $\mathbb{F}_2+v\mathbb{F}_2$, respectively. Then $d=d_G=\mathrm{min}\{d(\mathcal{C}_1),d(\mathcal{C}_2)\}$, where $d(\mathcal{C}_i)$ denotes the
minimum distance of binary codes $\mathcal{C}_i$ for $i=1,2$.
\end{proposition}

\begin{proof}
Since $\Theta$ is a distance preserving map, it follows that $d_G(\mathcal{C})=d(\Theta(\mathcal{C}))=d(\mathcal{C}_1 \otimes \mathcal{C}_2)=\mathrm{min}\{d(\mathcal{C}_1),d(\mathcal{C}_2)\}$. Thus, $d=d_G$.
\end{proof}

\begin{proposition}\label{prop2}
Let $\mathcal{C}=\langle \nu(x)\rangle$, where $\nu(x)=(r_1(x),(1+v)r_2(x)+vr_3(x))$ and $x^n-1=r_i(x)h_i(x)$ for all $i\in \{1,2,3\}$, be a cyclic code over $\mathcal{R}$. If $\mathcal{C}^{\perp_H}=\langle (\Bar{h}_1(x),v\Bar{h}_2(x)+(1+v)\Bar{h}_3(x))\rangle$, then $\tau(\mathcal{C})$ is a binary $[3n,k,d]$ code, where $k=3n-\mathrm{deg}(r_1(x))-\mathrm{deg}(r_2(x))-\mathrm{deg}(r_3(x))$ and $d=\mathrm{min}\{d(\langle r_1(x)\rangle),d(\langle r_2(x)\rangle),d(\langle r_3(x)\rangle)\}$.
\end{proposition}

\begin{proof}
The proof is the same as that of Lemma 2.6 in \cite{euclideansum}.
\end{proof}
\noindent In the next theorem, we present the generator polynomials and the corresponding dimensions of the Hermitian hulls and Hermitian sums associated with cyclic codes over the ring~$\mathcal{R}$.

\begin{theorem}
\label{th3} Let $\mathcal{C}=\langle \nu(x)\rangle$, where $\nu(x)=(r_1(x),(1+v)r_2(x)+r_3(x))$ and $\mathcal{C}^{\perp_H}=\langle \Bar{h}(x)\rangle$, where $\Bar{h}(x)=(\Bar{h}_1(x),v \Bar{h}_2(x)+(1+v)\Bar{h}_3(x))$ and $x^n-1=r_i(x)h_i(x)$ for all $i\in \{1,2,3\}$. Then
\begin{itemize}
\item[(i)] $\mathrm{Hull_{H}}(\mathcal{C})=\langle (\mathrm{lcm}(r_{1}(x),\Bar{h}_{1}(x)),(1+v)\mathrm{lcm}(r_{2}(x),\Bar{h}_{3}(x))+v\mathrm{lcm}(r_{3}(x),\Bar{h}_{2}(x)))\rangle $. Also,
\begin{equation*}
|\mathrm{Hull_{H}}(\mathcal{C})|=2^{3n-\mathrm{\deg }(\mathrm{lcm}(r_{1}(x),\Bar{h}_{1}(x)))-\mathrm{\deg }(\mathrm{lcm}(r_{2}(x),\Bar{h}_{3}(x)))-\mathrm{\deg }(\mathrm{lcm}(r_{3}(x),\Bar{h}_{2}(x)))}.
\end{equation*}

\item[(ii)] $\mathrm{Sum_{H}}(\mathcal{C})=\langle (\mathrm{\gcd }(r_{1}(x),\Bar{h}_{1}(x)),(1+v)\mathrm{\gcd }(r_{2}(x),\Bar{h}_{3}(x))+v\mathrm{\gcd }(r_{3}(x),\Bar{h}_{2}(x)))\rangle $. Also,
\begin{equation*}
|\mathrm{Sum_{H}}(\mathcal{C})|=2^{3n-\mathrm{\deg }(\mathrm{\gcd }(r_{1}(x),\Bar{h}_{1}(x)))-\mathrm{\deg }(\mathrm{\gcd }(r_{2}(x),\Bar{h}_{3}(x)))-\mathrm{\deg }(\mathrm{\gcd }(r_{3}(x),\Bar{h}_{2}(x)))}.
\end{equation*}
\end{itemize}
\end{theorem}

\begin{proof}
\begin{itemize}
\item[(i)] Assume that
\begin{equation*}
\mathcal{A}=\langle (c_{1}(x),(1+v)c_{2}(x)+vc_{3}(x))\rangle ,
\end{equation*}
where $c_{1}(x)=\mathrm{lcm}(r_{1}(x),\Bar{h}_{1}(x))$, $c_{2}(x)=\mathrm{lcm}(r_{2}(x),\Bar{h}_{3}(x))$ and $c_{3}(x)=\mathrm{lcm}(r_{3}(x),\Bar{h}_{2}(x)))$, is a cyclic code of length $n$ over $\mathcal{R}$. Then  there exist $a_{i}(x),$ $b_{i}(x)\in \mathbb{F}_{2}[x]$ and $i\in \{1,2,3\}$ such that
\begin{align*}
c_{1}(x)& =r_{1}(x)a_{1}(x)=\Bar{h}_{1}(x)b_{1}(x), \\
c_{2}(x)& =r_{2}(x)a_{2}(x)=\Bar{h}_{3}(x)b_{2}(x), \\
c_{3}(x)& =r_{3}(x)a_{3}(x)=\Bar{h}_{2}(x)b_{3}(x).
\end{align*}
So, $\mathcal{A}\subseteq \mathcal{C}$ and $\mathcal{A}\subseteq \mathcal{C}^{\perp }$ since
\begin{equation*}
(c_{1}(x),(1+v)c_{2}(x)+vc_{3}(x))=(a_{1}(x),(1+v)a_{2}(x)+va_{3}(x))\times
(r_{1}(x),(1+v)r_{2}(x)+vr_{3}(x)),
\end{equation*}
\begin{equation*}
(c_{1}(x),(1+v)c_{2}(x)+vc_{3}(x))=(b_{1}(x),(1+v)b_{2}(x)+vb_{3}(x))\times (\Bar{h}_{1}(x),v\Bar{h}_{2}(x)+(1+v)\Bar{h}_{3}(x)).
\end{equation*}
Therefore, $\mathcal{A}\subseteq \mathrm{Hull_{H}}(\mathcal{C})$. Clearly, $\mathrm{Hull_{H}}(\mathcal{C})$ is a cyclic code of length $n$ over $%
\mathcal{R}$. This implies that $\mathrm{Hull_{H}}(\mathcal{C})=\langle
(e_{1}(x),(1+v)e_{2}(x)+ve_{3}(x))\rangle $ for some $e_{1}(x),e_{2}(x),e_{3}(x)\in \mathbb{F}_{2}[x]$ such that $e_{i}(x)\mid
(x^{n}-1)$ for $i\in \{1,2,3\}$. Due to $\mathrm{Hull_{H}}(\mathcal{C}%
)\subset \mathcal{C}$,
\begin{align*}
(e_{1}(x),(1+v)e_{2}(x)+ve_{3}(x))&
=(s_{1}(x),(1+v)s_{2}(x)+vs_{3}(x))\times (r_{1}(x),(1+v)r_{2}(x)+vr_{3}(x))
\\
& =(s_{1}(x)r_{1}(x),(1+v)s_{2}(x)r_{2}(x)+vs_{3}(x)r_{3}(x))
\end{align*}
for some $s_{1}(x),$ $s_{2}(x),$ $s_{3}(x)\in \mathbb{F}_{2}[x]$. Then $%
r_{i}(x)\mid e_{i}(x)$ for $i\in \{1,2,3\}$. Similarly, $\mathrm{Hull_{H}}(%
\mathcal{C})\subset \mathcal{C}^{\perp }$ implies $\Bar{h}_{1}(x)\mid
e_{1}(x)$, $\Bar{h}_{2}(x)\mid e_{3}(x)$ and $\Bar{h}_{3}(x)\mid e_{2}(x)$.
Thus, $c_{i}(x)\mid e_{i}(x)$ and therefore there exists $p_{i}(x)\in
\mathbb{F}_{2}[x]$ such that $e_{i}(x)=c_{i}(x)p_{i}(x)$ for $i\in \{1,2,3\}$. Then
\begin{equation*}
(e_{1}(x),(1+v)e_{2}(x)+ve_{3}(x))=(p_{1}(x),(1+v)p_{2}(x)+vp_{3}(x))(c_{1}(x),(1+v)c_{2}(x)+vc_{3}(x)).
\end{equation*}
This means $\mathrm{Hull_{H}}(\mathcal{C})\subseteq \mathcal{A}$. Hence,
\begin{equation*}
\mathrm{Hull_{H}}(\mathcal{C})=\langle (\mathrm{lcm}(r_{1}(x),\Bar{h}_{1}(x)),(1+v) \mathrm{lcm}(r_{2}(x),\Bar{h}_{3}(x))+v\mathrm{\mathrm{lcm}}(r_{3}(x),\Bar{h}_{2}(x)))\rangle .
\end{equation*}
By Theorem \ref{7.1} part (i), we get
\begin{equation*}
|\mathrm{Hull_{H}}(\mathcal{C})|=2^{3n-\mathrm{\deg }(\mathrm{lcm}(r_{1}(x),\Bar{h}_{1}(x)))-\mathrm{\deg }(\mathrm{lcm}(r_{2}(x),\Bar{h}_{3}(x)))-\mathrm{\deg }(\mathrm{lcm}(r_{3}(x),\Bar{h}_{2}(x)))}.
\end{equation*}

\item[(ii)] The proof is similar to part (i).
\end{itemize}
\end{proof}

\section{QEC Codes from Hermitian Sums of Cyclic Codes over $\mathbb{F}_2\times(\mathbb{F}_2+v\mathbb{F}_2)$}
\label{sec4}
In this section, we obtain QEC codes by virtue of Hermitian hulls and sums of cyclic codes over $\mathcal{R}$.

\noindent An explicit method to get QEC codes from classical linear codes which is called quantum construction X of the Hermitian dual is given as follows.

\begin{theorem}[Quantum Construction X]
\cite{quantumxhermitian}\label{quantumxhermitian} Assume that $\mathcal{C}$ is an $[n,k]$-linear code over $\mathbb{F}_{p^2}$. Then there exists a quantum code over $\mathbb{F}_p$
with parameters $\llbracket n+e,2k-n,d\rrbracket $, where $d \geq \mathrm{min} \{d(\mathcal{C}%
),d(\mathrm{Sum_H}(\mathcal{C}))+1\}$ and $e = n-k-\mathrm{dim}(\mathrm{Hull_H}(\mathcal{C}))$.
\end{theorem}

\begin{theorem}
\label{th12} Let $\mathcal{C}=\langle
(r_{1}(x),(1+v)r_{2}(x)+r_{3}(x))\rangle $ and $x^{n}-1=r_{i}(x)h_{i}(x)$
for all $i\in \{1,2,3\}$. If
\begin{equation*}
\eta =\mathrm{deg}(\mathrm{lcm}(r_1(x),\Bar{h}_1(x)))+\mathrm{deg}(\mathrm{lcm}(r_2(x),\Bar{h}_3(x)))+\mathrm{deg}(\mathrm{lcm}(r_3(x),\Bar{h}_2(x)))
\end{equation*}
and
\begin{equation*}
\mu ={{\overset{3}{\underset{i=1}{{\sum }}}}}\mathrm{\deg }(r_{i}(x)),
\end{equation*}
then there exists a binary $\llbracket\eta +\mu -2n,n-2\mu ,d\geq \mathrm{\min }\{d(\tau (\mathcal{C})),d(\tau (\mathrm{Sum_{H}}(\mathcal{C})))+1\}\rrbracket$ QEC code, where
\begin{equation*}
d(\tau (\mathcal{C}))=\mathrm{\min }\{d(\langle r_{1}(x)\rangle ),d(\langle
r_{2}(x)\rangle ),d(\langle r_{3}(x)\rangle )\}
\end{equation*}
and
\begin{equation*}
d(\tau (\mathrm{Sum_{H}}(\mathcal{C})))=\mathrm{\min }\{d(\langle \mathrm{\gcd }(r_{1}(x),\Bar{h}_{1}(x))\rangle ),d(\langle \mathrm{\gcd }(r_{2}(x),\Bar{h}_{3}(x))\rangle ),d(\langle \mathrm{\gcd }(r_{3}(x),\Bar{h}_{2}(x))\rangle )\}.
\end{equation*}
\end{theorem}

\begin{proof}
By Theorem \ref{th3} part $(i)$,
\begin{equation*}
\mathrm{dim}(\mathrm{Hull_H}(\mathcal{C})))=3n-\mathrm{deg}(\mathrm{lcm}(r_1(x),\Bar{h}_1(x)))-\mathrm{deg}(\mathrm{lcm}(r_2(x),\Bar{h}_3(x)))-\mathrm{deg}(\mathrm{lcm}(r_3(x),\Bar{h}_2(x))).
\end{equation*}
From Theorem \ref{th3} part $(ii)$ and Lemma \ref{l2}, $d(\mathrm{Sum_{H}}(\tau (\mathcal{C})))=d(\tau (\mathrm{Sum_{H}}(\mathcal{C})))=\mathrm{\min } \{d(\langle \mathrm{\gcd }(r_{1}(x),\Bar{h}_{1}(x))\rangle ),d(\langle \mathrm{\gcd }(r_{2}(x),\Bar{h}_{3}(x))\rangle ),d(\langle \mathrm{\gcd }(r_{3}(x),\Bar{h}_{2}(x))\rangle )\}$. Moreover, $d(\tau (\mathcal{C}))=d(\mathcal{C})=\mathrm{\min }\{d(\langle r_{1}(x)\rangle ),d(\langle r_{2}(x)\rangle ),d(\langle r_{3}(x)\rangle )\}$. Hence, there exists a binary $\llbracket\eta +\mu -2n,n-2\mu ,d\geq \mathrm{\min }\{d(\tau (\mathcal{C})),d(\tau (\mathrm{Sum_{H}}(\mathcal{C})))+1\}\rrbracket$ QEC code by Theorem \ref{quantumxhermitian}.
\end{proof}

\begin{example}
The polynomial $x^{22}-1\in \mathbb{F}_{2}[x]$ is factorized as
\begin{equation*}
x^{22}-1=(x+1)^{2}\left(
x^{10}+x^{9}+x^{8}+x^{7}+x^{6}+x^{5}+x^{4}+x^{3}+x^{2}+x+1\right) ^{2}
\end{equation*}
Let $r_{1}(x)=x+1$ and $r_{2}(x)=r_{3}(x)=(x+1)^{2}$ and $\mathcal{C}=\langle (x+1)^{2},(1+v)(x+1)^{2}+(x+1)^{2})\rangle $. Then
\begin{eqnarray*}
\mathcal{C}^{\perp _{H}} &=&\langle ((x+1)\left(
x^{10}+x^{9}+x^{8}+x^{7}+x^{6}+x^{5}+x^{4}+x^{3}+x^{2}+x+1\right) ^{2}, \\
&&v\left( x^{10}+x^{9}+x^{8}+x^{7}+x^{6}+x^{5}+x^{4}+x^{3}+x^{2}+x+1\right)
\\
&&+(1+v)\left(
x^{10}+x^{9}+x^{8}+x^{7}+x^{6}+x^{5}+x^{4}+x^{3}+x^{2}+x+1\right) )\rangle .
\end{eqnarray*}
So,
\begin{equation*}
\mathrm{Sum_{H}}(\mathcal{C})=\langle ((x+1),(1+v)(1)+v(1))\rangle
\end{equation*}
and
\begin{eqnarray*}
\mathrm{Hull_{H}}(\mathcal{C}) &=&\langle ((x+1)\left(
x^{10}+x^{9}+x^{8}+x^{7}+x^{6}+x^{5}+x^{4}+x^{3}+x^{2}+x+1\right) ^{2}, \\
&&(1+v)(x^{22}-1)+v(x^{22}-1))\rangle
\end{eqnarray*}
Also,
\begin{equation*}
d(\tau (\mathcal{C}))=\mathrm{\min }\{d(\langle r_{1}(x)\rangle ),d(\langle
r_{2}(x)\rangle ),d(\langle r_{3}(x)\rangle )\} \\
=\mathrm{\min }\{2,2,2\} \\
=2
\end{equation*}
and
\begin{eqnarray*}
d(\tau (\mathrm{Sum_{H}}(\mathcal{C}))) &=&\mathrm{\min }\{d(\langle \mathrm{\gcd }(r_{1}(x),\Bar{h}_{1}(x))\rangle ), \\
&&d(\langle \mathrm{\gcd }(r_{2}(x),\Bar{h}_{3}(x))\rangle ),d(\langle
\mathrm{\gcd }(r_{3}(x),\Bar{h}_{2}(x))\rangle )\} \\
&=&\mathrm{\min }\{2,1,1\} \\
&=&1
\end{eqnarray*}
Thus, there exists a binary $\llbracket26,12,\geq 2\rrbracket$ QEC code by
Theorem \ref{th12}.
\end{example}
In Table \ref{TableTh10}, we tabulated some parameters of QECCs based on Theorem \ref{th12}. Most of the obtained codes are among the best known QECCs  according to \texttt{https://www.codetables.de/}( \cite{Grassl:codetables}) and some of them are optimal with respect to quantum singleton bound. All the computations are done using \verb"MAGMA" computer program \cite{bosma1997magma}.

\begin{table}[!ht]
\caption{Some binary QEC Codes from Theorem \ref{th12}}\label{TableTh10}
\begin{tabular}{lccc}
\toprule
$n$ & generators of $\mathcal{C}_{1}$ & generators of $\mathcal{C}_{2}$ & QECC
parameters \\
\midrule
$4$ & $1$ & $(1+v)(1) + v(x + 1)$ & $\llbracket 4,2,d\ge 2\rrbracket$   \\
$15$ & $1$ & $(1+v)(1)+v(x^{4}+x+1)$ & $\llbracket15,7,d\geq 3\rrbracket$ \\
$15$ & $1$ & $(1+v)(1)+v(x^{5}+x^{3}+x+1)$ & $\llbracket15,5,d\geq 4%
\rrbracket$ \\
$15$ & $1$ & $(1+v)(1)+v(x^{6}+x^{4}+x^{3}+x^{2}+1)$ & $\llbracket15,3,d\geq
4\rrbracket$ \\
$15$ & $x^{5}+x^{3}+x+1$ & $(1+v)(1)+v(1)$ & $\llbracket16,5,d\geq 4%
\rrbracket$ \\
$26$ & $1$ & $(1+v)(1)+v(1+x)$ & $\llbracket26,24,d\geq 2\rrbracket$ \\
$26$ & $1$ & $(1+v)(1)+v(x^{2}+1)$ & $\llbracket26,22,d\geq 2\rrbracket$ \\
$26$ & $1+x$ & $(1+v)(1)+v(x^{2}+1)$ & $\llbracket26,20,d\geq 2\rrbracket$
\\
$28$ & $1$ & $(1+v)(1)+v(1+x)$ & $\llbracket28,26,d\geq 2\rrbracket$
\\
$28$ & $1$ & $(1+v)(1)+v(x^{2}+1)$ & $\llbracket28,24,d\geq 2\rrbracket$ \\
$28$ & $1$ & $(1+v)(1)+v(x^{3}+x+1)$ & $\llbracket28,22,d\geq 2\rrbracket$
\\
$28$ & $1$ & $(1+v)(1)+v(x^{6}+x^{3}+x+1)$ & $\llbracket28,16,d\geq 4%
\rrbracket$ \\
$28$ & $x^{6}+x^{3}+x+1$ & $(1+v)(1)+v(1)$ & $\llbracket30,16,d\geq 3%
\rrbracket$ \\
$31$ & $1$ & $(1+v)(1)+v(1+x)$ & $\llbracket31,29,d= 1\rrbracket$ \\
$31$ & $1$ & $(1+v)(1)+v(x^{5}+x^{2}+1)$ & $\llbracket31,21,d\geq 3\rrbracket
$ \\
$31$ & $1$ & $(1+v)(1)+v(x^{6}+x^{2}+x+1)$ & $\llbracket31,19,d\geq 4%
\rrbracket$ \\
$31$ & $1+x$ & $(1+v)(1)+v(1+x)$ & $\llbracket32,27,d\geq 2\rrbracket$ \\
$31$ & $x^{6}+x^{2}+x+1$ & $(1+v)(1)+v(1)$ & $\llbracket32,19,d\geq 4%
\rrbracket$ \\
$31$ & $x^{11}+x^{6}+x^{5}+x^{2}+x+1$ & $(1+v)(1)+v(1)$ & $\llbracket%
32,9,d\geq 6\rrbracket$ \\
$36$ & $1$ & $(1+v)(1)+v(1+x)$ & $\llbracket36,34,d\geq 2\rrbracket$ \\
$36$ & $1$ & $(1+v)(1)+v(x^{2}+1)$ & $\llbracket36,32,d\geq 2\rrbracket$ \\
$36$ & $1$ & $(1+v)(1)+v(x^{3}+1)$ & $\llbracket36,30,d\geq 2\rrbracket$ \\
$36$ & $1$ & $(1+v)(1)+v(x^{9}+x^{8}+x^{7}+x^{5}+x^{4}+x^{2}+x+1)$ & $%
\llbracket36,18,d\geq 4\rrbracket$ \\
$48$ & $1$ & $(1+v)(1)+v(1+x)$ & $\llbracket48,46,d\geq 2\rrbracket$ \\
$48$ & $1$ & $(1+v)(1)+v(x^{2}+1)$ & $\llbracket48,44,d\geq 2\rrbracket$ \\
$48$ & $1$ & $(1+v)(1)+v(x^{3}+1)$ & $\llbracket48,42,d\geq 2\rrbracket$ \\
\botrule
\end{tabular}
\end{table}

\section{EAQEC codes from Hermitian Sums of Cyclic Codes over $\mathbb{F}_2 \times (\mathbb{F}_2+v\mathbb{F}_2)$}
\label{sec5}
In this section, we get EAQEC codes by using the following proposition.

\begin{proposition}\label{eaqec}\cite{eaqec}
Assume that $\mathcal{C}$ is a $q$-ary $[n,k,d]$-code. If $\mathcal{C}^{\perp}$ is its Euclidean dual code with parameters $[n,n - k,d^{\perp}]_q$, then there exist $q$-ary $\llbracket n,k -\mathrm{dim}(\mathrm{Hull}_{E}(\mathcal{C})),d;n-k-\mathrm{dim}(\mathrm{Hull}_{E}(\mathcal{C}))\rrbracket $ and $\llbracket n,n-k-\mathrm{dim}(\mathrm{Hull}_{E}(\mathcal{C})),d^\perp;k-\mathrm{dim}(\mathrm{Hull}_{E}(\mathcal{C}))\rrbracket $
 EAQEC codes.
\end{proposition}

\noindent Let $\mathcal{C}$ be a linear code. If $\mathrm{Hull}_{E}(\mathcal{C})=\{\mathbf{0}\}$, then $\mathcal{C}$ is named an LCD code. Similarly, it is termed as Hermitian LCD code when $\mathrm{Hull_H}(\mathcal{C})=\{\mathbf{0}\}$.

\begin{lemma}\cite{lcd}
Suppose that $\mathcal{C}$ is an $[n,k]$-code. If $G_{k\times n}$ is generator matrix of $\mathcal{C}$, then $\mathcal{C}$ is Hermitian LCD code if and only if $\mathrm{rk}(G\Bar{G}^T)=k$, where $\mathrm{rk}$ and $\Bar{G}^T$ stand for the rank of a matrix and the transpose of conjugate of $G$, respectively.
\end{lemma}

\begin{corollary}\label{corollary1}
Let $\mathcal{C}$ be a linear code of length $n$ and dimension $k$ over the ring~$\mathcal{R}$.
Then $\mathcal{C}$ is a Hermitian LCD code precisely when its Gray image $\tau(\mathcal{C})$ forms a binary LCD code with parameters~$[3n,\,k]$.
\end{corollary}

\begin{proof}
It is clear by Lemma \ref{l2}.
\end{proof}

\noindent A polynomial $s(x)$ of degree $t$ is termed as self-reciprocal if $s(x)=x^ts(\frac{1}{x})$.

\begin{proposition}\label{prop14}
Let $\mathcal{C} = \langle (r_1(x), (1+v)r_2(x) + r_3(x)) \rangle$ be a cyclic code of length~$n$ over the ring~$\mathcal{R}$,
where $x^n - 1 = r_i(x)h_i(x)$ in~$\mathbb{F}_2[x]$ and each $r_i(x)$ is self-reciprocal for $i = 1,2,3$.
If every monic irreducible factor of $r_i(x)$ occurs with the same multiplicity in both $r_i(x)$ and $x^n - 1$ for all $i \in \{1,2,3\}$,
then $\mathcal{C}$ is a Hermitian LCD code.
\end{proposition}

\begin{proof}
Similar with the proof of Proposition 4.6 in \cite{euclideansum}.
\end{proof}

\begin{theorem}\label{last_theorem}
Let $\mathcal{C}_{1}=\langle (r_{1}(x),(1+v)r_{2}(x)+vr_{3}(x))\rangle $ and
$\mathcal{C}_{2}=\langle (s_{1}(x),(1+v)s_{2}(x)+vs_{3}(x))\rangle $ be two
cyclic codes of length $n$ over the ring $\mathcal{R}$. Assume that the
polynomials $r_{i}(x)$ and $s_{i}(x)$ are self-reciprocal, and that each
monic irreducible factor of $r_{i}(x)$ and $s_{i}(x)$ appears with the same
multiplicity in both $r_{i}(x),s_{i}(x)$ and $x^{n}-1$ for all $i\in
\{1,2,3\}$. Define
\begin{equation*}
L=[\,\tau (\mathcal{C}_{1}),\,\tau (\mathcal{C}_{2})\,]B,
\end{equation*}
where
\begin{equation*}
B=
\begin{bmatrix}
1 & 1 & 0 & 1 \\
0 & 1 & 1 & 1
\end{bmatrix}
.
\end{equation*}
Then $L$ is a binary LCD code with parameters $[12n,\,k,\,d\geq \min
\{3d_{1},\,2d_{2}\}]_{2}$, where
\begin{equation*}
k=6n-\sum_{i=1}^{3}\deg (r_{i}(x))-\sum_{j=1}^{3}\deg (s_{j}(x)),
\end{equation*}
\begin{equation*}
d_{1}=\min \{d(\langle r_{1}(x)\rangle ),\,d(\langle r_{2}(x)\rangle
),\,d(\langle r_{3}(x)\rangle )\},
\end{equation*}
and
\begin{equation*}
d_{2}=\min \{d(\langle s_{1}(x)\rangle ),\,d(\langle s_{2}(x)\rangle
),\,d(\langle s_{3}(x)\rangle )\}.
\end{equation*}
Moreover, there exist binary EAQEC codes with parameters $\llbracket 12n,\,k,\,d\geq \min \{3d_{1},\,2d_{2}\};\,12n-k\rrbracket$.
\end{theorem}

\begin{proof}
The cyclic codes $\mathcal{C}_1$ and $\mathcal{C}_2$ are Hermitian LCD codes over $\mathcal{R}$ whose lengths are $n$ by Proposition \ref{prop14}. Also, $\tau(\mathcal{C}_1)$ and $\tau(\mathcal{C}_2)$ are LCD codes with the parameters $[3n,k_1,d_1]$ and $[3n,k_2,d_2]$, respectively, where
\begin{equation*}
   k_1 = {{\overset{3}{\underset{i=1} {{\displaystyle\sum}}}}}\mathrm{deg}(r_i(x)),
\end{equation*}
and
\begin{equation*}
   k_2 = {{\overset{3}{\underset{j=1} {{\displaystyle\sum}}}}}\mathrm{deg}(s_j(x)),
\end{equation*}
and
\begin{equation*}
d_1=\mathrm{min}\{d(\langle r_1(x)\rangle),d(\langle r_2(x)\rangle),d(\langle r_3(x)\rangle)\},
\end{equation*}
and
\begin{equation*}
d_2=\mathrm{min}\{d(\langle s_1(x)\rangle),d(\langle s_2(x)\rangle),d(\langle s_3(x)\rangle)\}.
\end{equation*}
from Corollary \ref{corollary1} and Proposition \ref{prop2}. The remaining part is similar with the Theorem 4.9 in \cite{euclideansum}. Moreover, by part (i) and Proposition \ref{eaqec}, there exists a binary $\llbracket 12n,k,d\geq \mathrm{min}\{3d_1,2d_2\};12n-k\rrbracket$ EAQEC code.
\end{proof}

\begin{example}
$\mathcal{C}_{1}=\langle (1+x),(1+v)(1+x)+v(1+x+x^2))\rangle $ and
$\mathcal{C}_{2}=\langle ((1+x+x^2),(1+v)(1+x+x^2)+v(1+x))\rangle $ are two
cyclic codes of length $9$ over the ring $\mathcal{R}$. In this case, the dimension of $\mathcal{C}$ is $54$ and $\mathrm{min}\{3d_1,2d_2\}=\mathrm{min}\{6,4\}=4$. Thus, there exists a binary $\llbracket 108,54,d\geq 4;54\rrbracket$ EAQEC code.
\end{example}

Using Theorem \ref{last_theorem}, we establish the existence of several new EAQEC codes for $n=3$ and $n=5$, as detailed in Tables \ref{n=3case} and \ref{n=5case}. These codes are not included in the EAQECC section of \texttt{https://www.codetables.de/} \cite{Grassl:codetables}. The parameters of even more new EAQEC codes, also derived from Theorem \ref{last_theorem} for other values of $n$, are tabulated in Table \ref{eaqectable}. All the computations are done using \verb"MAGMA" \cite{bosma1997magma}.

\begin{table}[!ht]
\caption{Some binary EAQEC Codes from Theorem \ref{last_theorem} in case of $n=3$}\label{n=3case}
 \begin{tabular}{@{}ll@{}}
 \toprule
 $r_i(x),s_i(x),1\leq i\leq 3$ & New EAQEC codes\\
 \midrule
 All of the polynomials are $1+x$. &  $\llbracket 36,12,d\geq 4;24\rrbracket_2$ \\
 One of $r_i(x)$ is $1+x$. The others are $1+x+x^2$. & $\llbracket 36,7,d\geq 4;29\rrbracket_2$ \\
 Two of $r_i(x)$ are $1+x$. The others are $1+x+x^2$. & $\llbracket 36,8,d\geq 4;28\rrbracket_2$ \\
 Three of $r_i(x)$ are $1+x$. The others are $1+x+x^2$. & $\llbracket 36,9,d\geq 4;27\rrbracket_2$ \\
 Four of $r_i(x)$ are $1+x$.The others are $1+x+x^2$. & $\llbracket 36,10,d\geq 4;26\rrbracket_2$ \\
 Five of $r_i(x)$ are $1+x$. The other is $1+x+x^2$. & $\llbracket 36,11,d\geq 4;25\rrbracket_2$ \\
  All of the polynomials are $1+x+x^2$. & $\llbracket 36,6,d\geq 6;30\rrbracket_2$ \\
 \botrule
 \end{tabular}
\end{table}

\begin{table}[!ht]
\caption{Some binary EAQEC Codes from Theorem \ref{last_theorem} in case of $n=5$}\label{n=5case}
 \begin{tabular}{@{}ll@{}}
 \toprule
 $r_i(x),s_i(x),1\leq i\leq 3$ & New EAQEC codes\\
 \midrule
 All of the polynomials are $1+x$. &  $\llbracket 60,24,d\geq 4;36\rrbracket_2$ \\
 One of $r_i(x)$ is $1+x$. The others are $1+x+x^2+x^3+x^4$. & $\llbracket 60,9,d\geq 2;51\rrbracket_2$ \\
 Two of $r_i(x)$ are $1+x$. The others are $1+x+x^2+x^3+x^4$. & $\llbracket 60,12,d\geq 4;48\rrbracket_2$ \\
 Three of $r_i(x)$ are $1+x$. The others are $1+x+x^2+x^3+x^4$. & $\llbracket 60,15,d\geq 4;45\rrbracket_2$ \\
 Four of $r_i(x)$ are $1+x$.The others are $1+x+x^2+x^3+x^4$. & $\llbracket 60,18,d\geq 4;42\rrbracket_2$ \\
 Five of $r_i(x)$ are $1+x$. The other is $1+x+x^2+x^3+x^4$. & $\llbracket 60,21,d\geq 4;39\rrbracket_2$ \\
  All of the polynomials are $1+x+x^2+x^3+x^4$. & $\llbracket 60,6,d\geq 10;54\rrbracket_2$ \\
 \botrule
 \end{tabular}
\end{table}

\begin{table}[!ht]
\caption{Some binary EAQEC Codes from Theorem \ref{last_theorem}}\label{eaqectable}
 \begin{tabular}{@{}ll@{}}
 \toprule
 $r_i(x),s_i(x),1\leq i\leq 3$ & New EAQEC codes\\
 \midrule
 $r_1(x)=r_2(x)=s_3(x)=1+x^2+x^3$, $s_1(x)=s_2(x)=r_3(x)=1+x$. &  $\llbracket 84,30,d\geq 4;64\rrbracket_2$ \\
 $r_1(x)=r_2(x)=s_3(x)=1+x$, $s_1(x)=s_2(x)=r_3(x)=1+x+x^2$. &  $\llbracket 108,54,d\geq 4;54\rrbracket_2$ \\
 $r_1(x)=s_1(x)=s_3(x)=1+x$, $s_2(x)=r_3(x)=1+x+x^2$, $r_2(x)=1+x+x^4$. &  $\llbracket 204,91,d\geq 10;113\rrbracket_2$ \\
 $r_1(x)=s_1(x)=s_3(x)=1+x+x^3$, $s_2(x)=r_3(x)=1+x$, $r_2(x)=1+x^2+x^3$. &  $\llbracket 252,112,d\geq 4;140\rrbracket_2$ \\

 \botrule
 \end{tabular}
\end{table}

\section{Conclusion} \label{sec6}
In this study, we first obtain QEC codes from the Hermitian sums and hulls of cyclic codes over the ring $\mathbb{F}_2\times(\mathbb{F}_2+v\mathbb{F}_2)$, via the Gray map $\tau$. Next, we examine the relationship between the Hermitian hulls of cyclic codes $\mathcal{C}$ over $\mathbb{F}_2\times(\mathbb{F}_2+v\mathbb{F}_2)$ and the Euclidean hulls of their Gray images
$\tau(\mathcal{C})$, which are linear codes of length $3n$ over $\mathbb{F}_2$. This relationship is then used to get binary QEC codes via the Hermitian dual-based quantum construction $X$. Using Hermitian sums and hulls, we obtained new QEC code parameters according to the Euclidean case. Finally, we propose a method to construct EAQEC codes by utilizing matrix product codes of LCD codes over the ring $\mathbb{F}_2\times(\mathbb{F}_2+v\mathbb{F}_2)$. We believe that the theory presented in this work could be generalized for primes $p>2$, which might yield new and good parameters for both QECCs and EAQECCs.

%\bmhead{Acknowledgements}

%The authors would like to thank the editors and the referees for their constructive comments, which have greatly contributed to improving the paper.

\section*{Declarations}

All authors declare that they have no conflict of interest. 

%%===========================================================================================%%
%% If you are submitting to one of the Nature Portfolio journals, using the eJP submission   %%
%% system, please include the references within the manuscript file itself. You may do this  %%
%% by copying the reference list from your .bbl file, paste it into the main manuscript .tex %%
%% file, and delete the associated \verb+\bibliography+ commands.                            %%
%%===========================================================================================%%

\bibliography{sn-bibliography}% common bib file
%% if required, the content of .bbl file can be included here once bbl is generated
%%\input sn-article.bbl

\end{document}
